% subsection: 対数型 \, $a_{n+1}=pa_n^q$

\subsection{\texorpdfstring{対数型 \, $a_{n+1}=pa_n^q$}{対数型}}\label{節：対数型}
  \begin{bookpoint}
    \textbf{積と冪をほどく}\par
    \bookterm{logarithm}{対数}を取る発想は指数の扱いを簡単にするためである.実数の\bookterm{logarithm}{対数}を使う前に,項や底の条件を確認する.
  \end{bookpoint}
  \noindent\textbf{【観察と発想】}\par
  \sectionref*{節：指数対数}の公式$\log_t(x^q)=q\log_t x$で,指数を係数にできる.
  積・商・累乗を和・差・定数倍へ変えるために,正値を確認して\bookterm{logarithm}{対数}を取ろう.

  \[
    a_{n+1}=pa_n^q
  \]
  ただし,$p>0,\,a_1>0$とする. このとき,すべての項が正となり,\bookterm{logarithm}{対数}を取ることができる.\footnote{$a_n>0$ならば$p a_n^q>0$であるから,\bookterm{induction}{数学的帰納法}によりすべての項が正となり,\bookterm{logarithm}{対数}を取ることができる.}
  まず,$p\ne1$の場合を考え,$p$を\bookterm{logarithm}{対数}の底に選ぶ.
  この式について,両辺の\bookterm{logarithm}{対数}を取ると
  \[
    \begin{aligned}
      \log_pa_{n+1}&=\log_pp+q\log_pa_n\\
                  &=q\log_pa_n+1.
    \end{aligned}
  \]
  $b_n=\log_pa_n$なら$b_{n+1}=qb_n+1$である.
  $q\ne1$ならば不動点を引いて定数項を消すことができる.
  $q=1$のときは$b_{n+1}=b_n+1$という等差型となり,$a_n=a_1p^{n-1}$を得る.
  また,$p=1$のときは$p$を\bookterm{logarithm}{対数}の底にできないが,底を$2$などに選べば
  \[
    b_n=\log_2a_n,\qquad b_{n+1}=qb_n
  \]
  となり,$q\ne0$ならば$a_n=a_1^{q^{n-1}}$を得る.
  $q=0$の場合は$a_1$を別に扱い,$n\ge2$では$a_n=1$である.
  底は正で$1$以外なら自由である.整数$m$で$p=k^m$と書けるなら,$k$を底にすると$\log_kp=m$となり,定数を整数にできる.
  例題(2)では商を差に変え,逆数の更新を符号反転として追う.

  \noindent\bookblocklink{example-logarithmic-1}{problem}{answer}{【例題】}\par\nobreak
  \begin{enumerate}[label=(\arabic*),parsep=3pt,format=\bookitemlink{example-logarithmic-1}{problem}{answer}]
    \item $\displaystyle a_1=2,\quad a_{n+1}=2a_n^2$
    \item $\displaystyle a_1=\frac{1}{2},\quad a_{n+1}=\frac{1}{a_n}$
    \item $\displaystyle a_1=1,\quad a_{n+1}=2^na_n$
    \item $\displaystyle a_1=5,\quad a_{n+1}=a_n^2-2a_n+2$
    \item $\displaystyle a_1=2,\quad a_{n+1}=\frac{a_n^2}{2^n}$
    \item $\displaystyle a_1=1,\quad a_{n+1}=(n+1)a_n$
  \end{enumerate}

  \noindent\bookblocklink{example-logarithmic-1}{hint}{problem}{【方針】}\par\nobreak
  \begin{enumerate}[label=(\arabic*),parsep=3pt,format=\bookitemlink{example-logarithmic-1}{hint}{problem}]
    \item $b_n=\log_2a_n$で$b_{n+1}=2b_n+1$に変え,定数項を消します.
    \item 逆数を\bookterm{logarithm}{対数}で符号反転に変えます.
    \item $f(n)=2^n$が指数関数なので,\bookterm{logarithm}{対数}を取ると和の形に帰着できます.
    \item \bookterm{square}{平方完成}で$a_{n+1}-1=(a_n-1)^2$を作ります.\bookterm{logarithm}{対数}を取る量は$a_n$ではなく$a_n-1$であり,その正値を先に確認します.
    \item $b_n=\log_2a_n$から一次式$n+1$を引くと等比型になります.
    \item 積を直接まとめる方法と,\bookterm{logarithm}{対数}で和に変える方法を比べましょう.最後には同じ積$n!$を計算します.
  \end{enumerate}

  \noindent\bookblocklink{example-logarithmic-1}{answer}{problem}{【解答】}\par\nobreak
  \begin{enumerate}[label=(\arabic*),parsep=3pt,format=\bookitemlink{example-logarithmic-1}{answer}{problem}]
    \item $\displaystyle a_n=2^{2^{n}-1}$
    \item $a_n=2^{(-1)^{n}}$
    \item $a_n=2^{\frac{n(n-1)}{2}}$
    \item $a_n=2^{2^n}+1$
    \item $\displaystyle a_n=2^{\,n+1-2^{n-1}}$
    \item $a_n=n!$
  \end{enumerate}

  \noindent\bookblocklink{example-logarithmic-1}{solution}{problem}{【解法】}\par\nobreak

  (1),(2),(3),(5),(6)は初項が正で,正の項から次も正になる.\bookterm{induction}{数学的帰納法}より全項が正である.
  (4)は$a_1-1=4>0$と$a_{n+1}-1=(a_n-1)^2$から,$a_n-1>0$を保つ.この量に\bookterm{logarithm}{対数}を取る.
  \begin{enumerate}[label=(\arabic*),parsep=3pt,format=\bookitemlink{example-logarithmic-1}{solution}{problem}]
    \item $b_n=\log_2a_n$と置くと,
          \[
            b_{n+1}=2b_n+1,\qquad b_1=1.
          \]
          $b_n+1$は初項$2$,公比$2$の等比型なので,
          \[
            b_n+1=2^n,\qquad a_n=2^{b_n}=2^{2^n-1}.
          \]
    \item $b_n=\log_2a_n$と置くと,逆数は
          \[
            b_{n+1}=-b_n,\qquad b_1=\log_2\frac12=-1
          \]
          に変わる.従って$b_n=(-1)^n$であり,$a_n=2^{(-1)^n}$となる.
    \item 自然\bookterm{logarithm}{対数}を取れば$\log a_{n+1}=\log a_n+n\log2$である.差を足すと,
          \[
            \log a_n=\sum_{k=1}^{n-1}k\log2=\frac{n(n-1)}2\log2.
          \]
          従って$a_n=2^{n(n-1)/2}$となる.
    \item \bookterm{square}{平方完成}で$a_{n+1}-1=(a_n-1)^2$となる.\bookterm{logarithm}{対数}を取ると,
          \[
            \log_2(a_{n+1}-1)=2\log_2(a_n-1).
          \]
          初項は$\log_2(a_1-1)=2$だから,$\log_2(a_n-1)=2^n$であり,$a_n=2^{2^n}+1$となる.
    \item $b_n=\log_2a_n$と置くと,$b_{n+1}=2b_n-n,\ b_1=1$である.
          一次式$n+1$を引けば,
          \[
            b_{n+1}-(n+2)=2\{b_n-(n+1)\}.
          \]
          $c_n=b_n-(n+1)$は$c_{n+1}=2c_n,\ c_1=-1$を満たすので,
          \[
            c_n=-2^{n-1},\quad b_n=n+1-2^{n-1}.
          \]
          戻すと$a_n=2^{\,n+1-2^{n-1}}$となる.
    \item 自然\bookterm{logarithm}{対数}を取れば$\log a_{n+1}=\log a_n+\log(n+1)$である.
          \[
            \log a_n=\log1+\log2+\cdots+\log n=\log n!.
          \]
          従って$a_n=n!$となる.直接積を取っても同じ結果が得られる.
  \end{enumerate}
  \noindent\textbf{【補足】}

  (3),(6)の$\log$は底が$e$の自然\bookterm{logarithm}{対数}である.
  変形後の真数が$0$や負なら,実数の\bookterm{logarithm}{対数}は取れない.例えば(4)で初項が$1$なら定数\bookterm{sequence}{数列}$a_n=1$となり,$\log(a_n-1)$を取らずに扱う.

  \noindent\bookblocklink{check-log-domain}{problem}{answer}{【確認問題】}\par\nobreak
  \begin{enumerate}[format=\bookitemlink{check-log-domain}{problem}{answer}]
    \item $a_{n+1}=a_n^2$で,初項が$2,0,-2$の場合を比較せよ.
          最初から実数の\bookterm{logarithm}{対数}を取れるのはどの場合か.取れない場合にも\bookterm{sequence}{数列}は定義されるか.
  \end{enumerate}
  \noindent\bookblocklink{check-log-domain}{answer}{problem}{【確認の解答】}\par\nobreak
  \begin{enumerate}[format=\bookitemlink{check-log-domain}{answer}{problem}]
    \item 初項$2$ならすべての項が正である.
          初項$0$なら零列,初項$-2$なら第2項以降が正である.
          いずれも元の\bookterm{recurrence}{漸化式}は定義される.\bookterm{logarithm}{対数}を取る条件と,元の\bookterm{sequence}{数列}が存在する条件は区別する.
  \end{enumerate}
